
% -------------------------------------------------------
% INSERT
% -------------------------------------------------------

\subsubsection{I2: Bulk import ocen}
\begin{figure}[H]
    \centering
    \includegraphics[width=0.8\textwidth]{charts/medium/chart_I2.png}
    \caption{I2 -- Bulk import ocen}
\end{figure}

Wzrost wolumenu ujawnia, jak bardzo \texttt{INSERT} jest operacją
liniową w liczbie wstawianych wierszy. PostgreSQL i MySQL skaczą
z 270/190~ms na 1.2/1.4~s -- niemal pięciokrotnie. MongoDB zachowuje
przewagę (461~ms vs 1.2--3.5~s u konkurencji), ale traci ją po
założeniu indeksów (1.1~s -- prawie tyle co PostgreSQL). Neo4j
płaci najwyższą cenę -- 3.5~s niezależnie od indeksów -- co wynika
z konieczności tworzenia osobnych węzłów \texttt{:Rating} dla każdej
oceny i wiązania ich relacjami do \texttt{:Profile} i \texttt{:Content}.

\subsubsection{I3: Dodanie serialu z pełnym drzewem}
\begin{figure}[H]
    \centering
    \includegraphics[width=0.8\textwidth]{charts/medium/chart_I3.png}
    \caption{I3 -- Dodanie serialu z pełnym drzewem}
\end{figure}

Rozkład sił identyczny jak na \emph{small} -- dowód, że ten scenariusz
jest praktycznie \emph{niezależny od wolumenu danych} (wstawiamy
zawsze ten sam ,,jeden serial''). MongoDB \emph{0.7~ms}, pozostali
15--30~ms. Niezmienność wyniku potwierdza, że koszt \texttt{INSERT}-a
pojedynczej encji nie zależy od rozmiaru tabeli docelowej -- liczy
się tylko narzut samej operacji zapisu i utrzymania indeksów na tym
konkretnym wierszu.

% -------------------------------------------------------
% SELECT
% -------------------------------------------------------

\subsubsection{S2: Filtrowanie kolaboratywne}
\begin{figure}[H]
    \centering
    \includegraphics[width=0.8\textwidth]{charts/medium/chart_S2.png}
    \caption{S2 -- Filtrowanie kolaboratywne}
\end{figure}

Indeksy zaczynają działać spektakularnie. Z indeksami \emph{wszystkie}
silniki schodzą do 1.6--5.5~ms -- a bez nich PostgreSQL i MongoDB
mierzą się z 40--97~ms. To pierwszy moment, w którym hipoteza pracy
ma jednoznacznie potwierdzającą postać: indeks redukuje czas SELECT-a
o \emph{rząd wielkości lub więcej}. Neo4j najszybszy bez indeksów
(8~ms) -- struktura grafu zaczyna się tu opłacać.

\subsubsection{S3: TOP 100 według liczby ocen w okresie}
\begin{figure}[H]
    \centering
    \includegraphics[width=0.8\textwidth]{charts/medium/chart_S3.png}
    \caption{S3 -- TOP 100 według liczby ocen w okresie}
\end{figure}

Trend ze \emph{small} się utrzymuje -- indeks niewiele pomaga, bo
filtr ,,od czerwca'' pokrywa za szeroki podzbiór. PostgreSQL nadal
najszybszy (15--18~ms), Neo4j skacze do 100~ms ze względu na koszt
przechodzenia po wszystkich krawędziach \texttt{:WATCHED}.
MySQL stabilnie 60~ms. Brak istotnej różnicy między wariantami
indeksowymi w żadnym z silników.

\subsubsection{S4: Wyszukiwanie pełnotekstowe po tytule}
\begin{figure}[H]
    \centering
    \includegraphics[width=0.8\textwidth]{charts/medium/chart_S4.png}
    \caption{S4 -- Wyszukiwanie pełnotekstowe po tytule}
\end{figure}

Tu dopiero widać prawdziwą wartość indeksu pełnotekstowego.
Bez indeksów Neo4j wykonuje \texttt{AllNodesScan} na 5\,000 węzłach --
30~ms. Z indeksem (\texttt{db.index.fulltext.queryNodes}) schodzi do
4~ms. PostgreSQL z trygramowym GIN-em: 0.9~ms (poprawa o 4×).
MongoDB z indeksem \texttt{text}: 1.2~ms (poprawa o 4×). MySQL --
jak pokazaliśmy w 11.3 -- nie aktywuje indeksu \texttt{FULLTEXT}
dla operatora \texttt{LIKE '\%term\%'}, więc obie wartości są
porównywalne (2.8 / 1.3~ms).

% -------------------------------------------------------
% DELETE
% -------------------------------------------------------

\subsubsection{D1: Usunięcie treści z kaskadą}
\begin{figure}[H]
    \centering
    \includegraphics[width=0.8\textwidth]{charts/medium/chart_D1.png}
    \caption{D1 -- Usunięcie treści z kaskadą}
\end{figure}

Skala intensyfikuje obraz znany ze \emph{small}: PostgreSQL bez
indeksu 24~ms, z indeksem 1.3~ms (18-krotne przyspieszenie). MongoDB
bez indeksu 136~ms, z indeksem 57~ms (2.4-krotnie). MySQL i Neo4j
działają stabilnie 6 i 9~ms odpowiednio. Indeksy na FK są tu
\emph{kluczowe} -- bez nich kaskadowe usunięcie staje się sumą
pełnych skanów wszystkich tabel zależnych.

\subsubsection{D4: Usunięcie pozycji z listy}
\begin{figure}[H]
    \centering
    \includegraphics[width=0.8\textwidth]{charts/medium/chart_D4.png}
    \caption{D4 -- Usunięcie pozycji z listy}
\end{figure}

Trend liniowy MongoDB w wariancie bez indeksu się potwierdza --
20~ms na \emph{small}, 55~ms na \emph{medium}. Z indeksem to wciąż
0.7~ms, czyli \emph{80-krotna} różnica. PostgreSQL i MySQL stabilne
poniżej 6~ms. Ten prosty scenariusz najwyraźniej obrazuje regułę:
\texttt{deleteOne} bez indeksu jest \texttt{COLLSCAN}-em, koniec
historii.

\subsubsection{D6: Masowe usunięcie nieaktywnych użytkowników}
\begin{figure}[H]
    \centering
    \includegraphics[width=0.8\textwidth]{charts/medium/chart_D6.png}
    \caption{D6 -- Masowe usunięcie nieaktywnych użytkowników}
\end{figure}

Zaczyna być \emph{dramatycznie}. PostgreSQL bez indeksu na
\texttt{users.status} -- \textbf{1.4~sekundy}, z indeksem 3.3~ms.
To 425-krotne przyspieszenie -- jeden z najmocniejszych wyników
w pracy. MongoDB i Neo4j odpowiednio 2/4~ms i 28/10~ms -- też
zyskują, ale bez tej dramaturgii. MySQL stabilnie 8~ms. Na
\emph{large} ten kontrast pokaże się jeszcze mocniej -- PostgreSQL
przekroczy 30~sekund.

% -------------------------------------------------------
% UPDATE
% -------------------------------------------------------

\subsubsection{U2: Przeliczenie średniej oceny}
\begin{figure}[H]
    \centering
    \includegraphics[width=0.8\textwidth]{charts/medium/chart_U2.png}
    \caption{U2 -- Przeliczenie średniej oceny}
\end{figure}

PostgreSQL: 8.8 → 1.5~ms (5×). MongoDB: 30 → 1.6~ms (19×). MySQL
stabilnie 6~ms, Neo4j 4~ms. Pojedynczy \texttt{UPDATE} po znalezieniu
wiersza z agregatu, wzorzec znany z \emph{small} -- tylko skala
narzutu wzrasta tam, gdzie indeksu brakuje. Indeks robi tu robotę
oczekiwaną (przyspiesza punktowy filtr), bez negatywnego wpływu
na samą operację zapisu, bo modyfikujemy jeden wiersz.

\subsubsection{U3: Masowa zmiana planu subskrypcji}
\begin{figure}[H]
    \centering
    \includegraphics[width=0.8\textwidth]{charts/medium/chart_U3.png}
    \caption{U3 -- Masowa zmiana planu subskrypcji}
\end{figure}

Koszt utrzymania indeksów przy modyfikacji wierszy zaczyna być
zauważalny. PostgreSQL bez indeksu 30~ms,
z indeksem 65~ms (ponad dwukrotne pogorszenie). MySQL 124 → 179~ms,
Mongo 107 → 104~ms (tu indeks praktycznie neutralny). Neo4j
zachowuje się odwrotnie -- 194 → 158~ms (lekkie polepszenie),
prawdopodobnie dlatego że \texttt{plan\_name} nie jest indeksowany
w grafie. Im więcej zaktualizowanych wierszy i im więcej indeksów
na zmienianej kolumnie, tym wyższa kara czasowa -- to jest
podręcznikowy efekt utrzymania indeksu przy modyfikacji.
